% DFFRNT AI Assistant — Developer Guide
% Overleaf-compatible (pdfLaTeX). Styling matches the DFFRNT report theme:
% Palatino body, garnet (uOttawa-inspired) accent, markdown-style callouts,
% themed tables.
\documentclass[11pt,letterpaper]{article}

%------------------------------------------------------------------ encoding --
\usepackage[T1]{fontenc}
\usepackage[utf8]{inputenc}

% Silence the harmless "Command \showhyphens has changed" notice microtype
% emits under recent LaTeX kernels.
\usepackage{silence}
\WarningFilter{latex}{Command}

%--------------------------------------------------------------------- fonts --
\usepackage{mathpazo}                 % Palatino body + matching math
\usepackage[scaled=0.92]{helvet}      % Helvetica-like sans (headings)
\usepackage{microtype}                % refined justification & spacing
\linespread{1.08}
\renewcommand{\sfdefault}{\rmdefault} % unify on Palatino, matching the report

%-------------------------------------------------------------------- layout --
\usepackage[letterpaper,margin=1in,headheight=22pt,headsep=18pt,footskip=28pt]{geometry}
\usepackage{tocloft}                  % load BEFORE parskip
\usepackage{parskip}
\usepackage{enumitem}
\setlist{itemsep=2pt,topsep=4pt,parsep=0pt}
\usepackage{environ}

%--------------------------------------------------------------------- colour --
\usepackage[table]{xcolor}
\definecolor{accent}{HTML}{6A24B5}    % DFFRNT brand purple
\definecolor{accentdk}{HTML}{4E1A86}  % deeper purple (headings / links)
\definecolor{brandlime}{HTML}{CFEF3C} % chartreuse accent
\definecolor{pagebg}{HTML}{E7E2F0}    % lavender title-page background
\definecolor{icnk}{HTML}{171634}      % dark navy icon tile
\definecolor{blocklt}{HTML}{D9CEEE}   % light lavender text on the purple block
\definecolor{blocklt2}{HTML}{C7B9E6}  % lighter lavender (block small print)
\definecolor{ink}{HTML}{1F2933}       % near-black body ink
\definecolor{inkmute}{HTML}{56616C}   % muted grey for metadata
\definecolor{rulegrey}{HTML}{C7CDD4}  % hairline rules
\definecolor{briefbg}{HTML}{F1ECF8}   % light lavender callout / code background

% Admonition palette (markdown-style callouts)
\definecolor{noteblue}{HTML}{2B6CB0}
\definecolor{notebg}{HTML}{EEF4FB}
\definecolor{tipgreen}{HTML}{2F7A55}
\definecolor{tipbg}{HTML}{EDF6F0}
\definecolor{warnred}{HTML}{B02A37}
\definecolor{warnbg}{HTML}{FBEEEF}

%-------------------------------------------------------------------- tables ---
\usepackage{array}
\usepackage{booktabs}
\usepackage{tabularx}
\usepackage{xltabular}   % page-breakable tabularx (loads longtable + tabularx)
\definecolor{tablehead}{HTML}{6A24B5}   % header fill (= accent)
\definecolor{tableshade}{HTML}{F1ECF8}  % alternating-row shade

\newcolumntype{L}{>{\raggedright\arraybackslash}X}
\newcolumntype{C}{>{\centering\arraybackslash}X}
\newcommand{\thd}[1]{{\sffamily\bfseries\color{white}#1}}

% ptable: styled non-breaking table body. Argument = column spec.
\NewEnviron{ptable}[1]{%
  \par\vspace{4pt}\footnotesize
  \setlength{\extrarowheight}{3pt}%
  \arrayrulecolor{rulegrey}%
  \rowcolors{2}{tableshade}{white}%
  \begin{tabularx}{\linewidth}{#1}
  \rowcolor{tablehead}\BODY
  \end{tabularx}%
  \arrayrulecolor{black}\par\vspace{6pt}%
}

% pltable: page-breaking sibling (header declared with \endhead in the body).
\NewEnviron{pltable}[1]{%
  \footnotesize
  \setlength{\extrarowheight}{3pt}%
  \arrayrulecolor{rulegrey}%
  \rowcolors{2}{tableshade}{white}%
  \begin{xltabular}{\linewidth}{#1}
  \BODY
  \end{xltabular}%
  \arrayrulecolor{black}%
}

%------------------------------------------------------------- code listings ---
\usepackage{listings}
\lstset{
  basicstyle=\ttfamily\small,
  keywordstyle=\color{accentdk}\bfseries,
  commentstyle=\color{inkmute}\itshape,
  stringstyle=\color{accent},
  frame=single,
  rulecolor=\color{rulegrey}, framesep=6pt,
  showstringspaces=false, breaklines=true, tabsize=2,
  backgroundcolor=\color{briefbg!40},
  xleftmargin=4pt,
}
% Directory trees / diagrams: smaller, no frame colouring changes needed
\lstdefinestyle{tree}{basicstyle=\ttfamily\footnotesize}

%---------------------------------------------------- section / TOC styling ---
\usepackage{titlesec}
\titleformat{\section}
  {\sffamily\LARGE\bfseries\color{accentdk}}
  {\makebox[1.6em][l]{\thesection}}{0em}{}
  [\vspace{3pt}{\color{accent}\titlerule[1pt]}]
\titleformat{\subsection}
  {\sffamily\Large\bfseries\color{ink}}
  {\makebox[2.0em][l]{\thesubsection}}{0em}{}
\titleformat{\subsubsection}
  {\sffamily\large\bfseries\color{ink}}
  {\thesubsubsection}{0.6em}{}
% FAQ questions: garnet, run-in
\titleformat{\paragraph}[runin]
  {\sffamily\bfseries\color{accentdk}}{}{0em}{}
\titlespacing*{\section}      {0pt}{0pt}{14pt}
\titlespacing*{\subsection}   {0pt}{16pt}{6pt}
\titlespacing*{\subsubsection}{0pt}{12pt}{4pt}
\titlespacing*{\paragraph}    {0pt}{8pt}{0.7em}

\renewcommand{\thesection}{\arabic{section}}
\renewcommand*\contentsname{Table of Contents}
\renewcommand{\cfttoctitlefont}{\sffamily\LARGE\bfseries\color{accentdk}}
\renewcommand{\cftsecfont}{\sffamily\bfseries\color{ink}}
\renewcommand{\cftsecpagefont}{\sffamily\bfseries\color{ink}}
\renewcommand{\cftsubsecfont}{\color{inkmute}}
\renewcommand{\cftsubsecpagefont}{\color{inkmute}}
\setlength{\cftbeforesecskip}{6pt}
\renewcommand{\cftsecleader}{\cftdotfill{\cftdotsep}}

%-------------------------------------------------------- header / footer -----
\usepackage{fancyhdr}
\pagestyle{fancy}
\fancyhf{}
\renewcommand{\headrulewidth}{0.4pt}
\renewcommand{\headrule}{\hbox to\headwidth{\color{rulegrey}\leaders\hrule height \headrulewidth\hfill}}
\renewcommand{\footrulewidth}{0pt}
\renewcommand{\sectionmark}[1]{\markboth{\thesection\quad #1}{}}
\fancyhead[L]{\footnotesize\sffamily\color{inkmute}Developer Guide}
\fancyhead[R]{\footnotesize\sffamily\color{inkmute}\nouppercase{\leftmark}}
\fancyfoot[R]{\footnotesize\sffamily\color{inkmute}Page \thepage}
\fancyfoot[L]{\footnotesize\sffamily\color{inkmute}DFFRNT Air-Gapped AI Assistant}
\fancypagestyle{plain}{%
  \fancyhf{}%
  \renewcommand{\headrulewidth}{0pt}%
  \fancyfoot[R]{\footnotesize\sffamily\color{inkmute}Page \thepage}%
  \fancyfoot[L]{\footnotesize\sffamily\color{inkmute}DFFRNT Air-Gapped AI Assistant}%
}

%------------------------------------------------------- hyperlinks (late) ----
\usepackage[colorlinks=true,linkcolor=accentdk,urlcolor=accent,
            citecolor=accentdk,linktoc=all]{hyperref}
\usepackage{url}
\urlstyle{sf}

%------------------------------------------------ markdown-style callouts ------
\usepackage[most]{tcolorbox}
\usetikzlibrary{shapes.geometric}   % 4-point "sparkle" star on the title page
\usepackage{graphicx}
\graphicspath{{./}{doc/}}
% Lime-tile icon. Uses the real DFFRNT PNG logo when `dffrnt-icon.png` is present
% beside this file; otherwise falls back to a vector sparkle so the page always
% renders. Drop the PNG in to get the true logo; delete this macro's fallback
% branch (leaving just the circle) to have empty tiles instead.
\newcommand{\dffrntlogo}[1]{%
  \IfFileExists{dffrnt-icon.png}
    {\node at (#1) {\includegraphics[width=15mm]{dffrnt-icon.png}};}
    {\node[icontile] at (#1) {};
     \node[sparkle]  at (#1) {};
     \node[star, star points=4, star point ratio=0.34, fill=brandlime,
           minimum size=2.6mm, inner sep=0] at ([xshift=-3.9mm,yshift=-3.4mm]#1) {};}%
}
\newtcolorbox{gnote}{
  enhanced, breakable,
  colback=notebg, colframe=noteblue,
  boxrule=0pt, leftrule=3.5pt, sharp corners,
  left=12pt, right=12pt, top=8pt, bottom=8pt,
  before skip=10pt, after skip=12pt,
  fontupper=\small\color{ink},
  before upper={{\sffamily\bfseries\color{noteblue}Note}\par\smallskip},
}
\newtcolorbox{gtip}{
  enhanced, breakable,
  colback=tipbg, colframe=tipgreen,
  boxrule=0pt, leftrule=3.5pt, sharp corners,
  left=12pt, right=12pt, top=8pt, bottom=8pt,
  before skip=10pt, after skip=12pt,
  fontupper=\small\color{ink},
  before upper={{\sffamily\bfseries\color{tipgreen}Tip}\par\smallskip},
}
\newtcolorbox{gwarn}{
  enhanced, breakable,
  colback=warnbg, colframe=warnred,
  boxrule=0pt, leftrule=3.5pt, sharp corners,
  left=12pt, right=12pt, top=8pt, bottom=8pt,
  before skip=10pt, after skip=12pt,
  fontupper=\small\color{ink},
  before upper={{\sffamily\bfseries\color{warnred}Warning}\par\smallskip},
}

\setcounter{tocdepth}{2}
\setcounter{secnumdepth}{3}

% Start each major section a little lower.
\let\oldsection\section
\renewcommand{\section}{\vspace{1em}\oldsection}

%==============================================================================
\begin{document}
%==============================================================================

%--------------------------------------------------------------- TITLE PAGE ----
\begin{titlepage}
\thispagestyle{empty}
\begin{tikzpicture}[remember picture, overlay,
    icontile/.style={fill=icnk, rounded corners=2mm, minimum size=15mm, inner sep=0},
    sparkle/.style={star, star points=4, star point ratio=0.36, fill=brandlime,
                    minimum size=8mm, inner sep=0}]
  % page background
  \fill[pagebg] (current page.north west) rectangle (current page.south east);

  % --- decorative circle cluster (top) ---
  \fill[accent]    ([xshift=83mm, yshift=-40mm]current page.north west) circle (17mm);
  \fill[accent]    ([xshift=167mm,yshift=-43mm]current page.north west) circle (17mm);
  \fill[accent]    ([xshift=131mm,yshift=-85mm]current page.north west) circle (15mm);
  \fill[accent]    ([xshift=170mm,yshift=-122mm]current page.north west) circle (14mm);
  \coordinate (limeA) at ([xshift=128mm,yshift=-38mm]current page.north west);
  \coordinate (limeB) at ([xshift=172mm,yshift=-86mm]current page.north west);
  \fill[brandlime] (limeA) circle (17mm);
  \fill[brandlime] (limeB) circle (17mm);
  \dffrntlogo{limeA}
  \dffrntlogo{limeB}

  % --- product title + tagline ---
  \node[anchor=west, text=accentdk,
        font=\fontfamily{phv}\bfseries\fontsize{40}{46}\selectfont]
    at ([xshift=26mm,yshift=-152mm]current page.north west) {DFFRNT};
  \node[anchor=west, text=accentdk,
        font=\fontfamily{phv}\bfseries\fontsize{40}{46}\selectfont]
    at ([xshift=26mm,yshift=-168mm]current page.north west) {AI Assistant};
  \node[anchor=north west, text=accent, font=\fontfamily{phv}\selectfont\large]
    at ([xshift=26mm,yshift=-180mm]current page.north west)
    {\parbox{155mm}{A fully air-gapped, private knowledge and content AI assistant}};

  % --- guide block ---
  \fill[accent, rounded corners=5mm]
    ([xshift=20mm,yshift=-190mm]current page.north west)
    rectangle ([xshift=196mm,yshift=-258mm]current page.north west);
  \node[anchor=west, text=brandlime,
        font=\fontfamily{phv}\bfseries\fontsize{42}{46}\selectfont]
    at ([xshift=32mm,yshift=-212mm]current page.north west) {DEVELOPER GUIDE};
  \node[anchor=north west, text=blocklt, font=\fontfamily{phv}\selectfont\large]
    at ([xshift=32mm,yshift=-222mm]current page.north west)
    {\parbox{152mm}{Architecture, source tour, and extension guidance for
     future development teams}};
  \node[anchor=north west, text=blocklt2, font=\fontfamily{phv}\selectfont\footnotesize]
    at ([xshift=32mm,yshift=-234mm]current page.north west)
    {\parbox{152mm}{The technical companion to the Installer and User Guides:
     how the RAG pipeline is built, why the stack looks the way it does, how to
     set up a development environment, and where to put new capabilities.}};
  \node[anchor=west, text=brandlime, font=\fontfamily{phv}\bfseries\selectfont\small]
    at ([xshift=32mm,yshift=-252mm]current page.north west) {\today};
\end{tikzpicture}
\end{titlepage}

%---------------------------------------------------------- FRONT MATTER -------
\pagenumbering{roman}
\setcounter{page}{1}
\pagestyle{fancy}

{\hypersetup{linkcolor=ink}\tableofcontents}
\clearpage

%------------------------------------------------------------- MAIN MATTER ------
\pagenumbering{arabic}
\setcounter{page}{1}

This guide is for developers --- in particular future grad-student teams ---
picking up the DFFRNT Air-Gapped AI Assistant to extend its capabilities. It
assumes general software-engineering background but \textbf{zero familiarity
with this codebase}. Read it top to bottom once; afterwards the per-module
docstrings and the two companion documents (Installer Guide, User Guide)
fill in the operational details.

\section{Architecture and tech stack}

\subsection{What the system is}

A fully offline Retrieval-Augmented Generation (RAG) assistant. Company
documents are ingested into a local vector database; questions are answered
by a locally served LLM, grounded in retrieved document chunks, with inline
\texttt{[n]} citations. The defining constraint is the \textbf{air gap}:
after deployment, nothing may require internet access --- which drives
almost every technology choice below.

\subsection{The pipeline}

\begin{lstlisting}[style=tree]
              +---------- INGESTION -----------+
Documents --> loaders --> chunker --> embed (bge-m3) --> Qdrant
(PDF, DOCX,..) (text +    (512-char    via Ollama       (chunks + payload
                sections)  recursive)                    + per-doc summaries)

              +---------- QUERY ---------------+
Question --> embed --> Qdrant search --> context assembly --> LLM (qwen3)
                       (top-k + margin cut         |            via Ollama
                        + aggregate grouping)      v
                                  prompt: system + <documents> + history
                                                   |
                                                   v
                                streamed answer + [n] citations + sources
\end{lstlisting}

\emph{Ingestion:} a format-specific \textbf{loader} produces clean text with
section metadata (page numbers, slide indexes, headings); the
\textbf{chunker} splits it (recursive strategy, 512 chars, 64 overlap, 128
floor); each chunk is \textbf{embedded} by \texttt{bge-m3} through Ollama
and stored in \textbf{Qdrant} with a rich payload (see
\texttt{schema.py}). A per-document \textbf{summary} is also generated and
stored in a second collection --- it routes aggregate queries and briefs the
prompt.

\emph{Query:} the question is embedded, Qdrant returns the \texttt{top\_k}
nearest chunks, a \textbf{score-margin cut} drops weak hits, and --- for
aggregate questions (``compare all candidates'') or tag-scoped queries ---
a \textbf{grouped search} additionally fetches the best chunk per document
so roll-ups cover every relevant file. The chunks are rendered into a
\texttt{<documents>} block, combined with the system prompt and recent
conversation history, and streamed through the LLM. The UI renders the
stream (optionally including the model's reasoning), maps \texttt{[n]}
citations to sources, and persists the conversation.

\subsection{Runtime topology}

Three services, wired by Docker Compose
(\texttt{deploy/docker-compose.yml}):

\begin{ptable}{>{\ttfamily}p{1.6cm} p{4.4cm} L >{\ttfamily}p{1.2cm}}
\thd{Service} & \thd{Image} & \thd{Role} & \thd{Port}\\
qdrant & qdrant/qdrant (pinned) & Vector store + document/tag/conversation
persistence & 6333\\
ollama & ollama/ollama (pinned) & Serves the LLM and the embedder & 11434\\
api & dffrnt-assistant:latest (built from deploy/Dockerfile) & FastAPI app +
static UI & 8000\\
\end{ptable}

In \textbf{development} the API usually runs on the host
(debugger-friendly) against containerized Qdrant + Ollama. In
\textbf{production} all three are containers (\texttt{prod} compose
profile).

There are two deployment \emph{pathways}, abstracted behind an interface
(\S\ref{sec:structure}, \texttt{installer/backends/}):

\begin{itemize}
  \item \textbf{CONTAINER} (Linux, Windows) --- everything as above.
  \item \textbf{PORTABLE} (macOS) --- Docker's Linux VM has no Metal GPU
        passthrough, so Ollama runs as a \textbf{native host process}
        (Metal is automatic on Apple Silicon) while Qdrant + API run in a
        bundled portable container runtime (colima + lima + static docker
        CLI).
\end{itemize}

\subsection{Tech stack}

\begin{pltable}{p{3.1cm} p{4.6cm} L}
\rowcolor{tablehead}\thd{Layer} & \thd{Technology} & \thd{Notes}\\
\endfirsthead
\rowcolor{tablehead}\thd{Layer} & \thd{Technology} & \thd{Notes}\\
\endhead
Language & Python $\geq$ 3.14 & stdlib \texttt{tomllib} for config, modern
typing\\
Package manager & uv & lockfile (\texttt{uv.lock}), dependency groups\\
API & FastAPI + uvicorn & thin HTTP layer; logic lives in services\\
Vector DB & Qdrant (\texttt{qdrant-client}) & also (ab)used as a plain
document store\\
LLM serving & Ollama & \texttt{qwen3} family (LLM) + \texttt{bge-m3}
(embeddings)\\
Ollama client & \textbf{stdlib \texttt{urllib}} (\texttt{ollama.py}) &
deliberately no SDK --- see FAQ\\
Document parsing & PyMuPDF, python-docx, python-pptx, openpyxl, stdlib
\texttt{csv} & one library per format; \textbf{no LangChain /
unstructured}\\
Answer $\rightarrow$ PDF & mistune (MD$\rightarrow$HTML) + PyMuPDF Story
(HTML$\rightarrow$PDF) & no extra PDF engine\\
Web UI & Vanilla HTML/CSS/JS, native ES modules & \textbf{no npm, no
framework, no build step}\\
Manager UI & stdlib \texttt{http.server} + static HTML/JS & the browser is
the GUI toolkit\\
Packaging & Docker + PyInstaller (\texttt{dffrnt-manager}) & macOS ships
source + vendored CPython instead\\
Tests & pytest (3 suites: unit / eval / integration) & plus a synthetic
eval corpus generator\\
\end{pltable}

The unifying philosophy: \textbf{minimum dependency surface}. Every
dependency must earn its place, because everything has to ship inside an
offline bundle and keep working with no ability to \texttt{pip install} a
fix on the target.

\section{Source code structure}
\label{sec:structure}

\begin{lstlisting}[style=tree]
dffrnt-airgapped-ai-assistant/
  config.toml             live config -- ALSO frozen into bundles at package time
  config.toml.example     annotated template (fallback for a bare checkout)
  conftest.py             pytest path setup + --reset-corpus option
  pyproject.toml          deps, dependency groups (dev/package), pytest config
  uv.lock                 locked dependency versions (committed)

  dffrnt_assistant/       === THE APPLICATION ===
    config.py             Settings dataclass: every tunable, env>TOML>default
    schema.py             THE shared payload contract (chunk fields, helpers)
    ollama.py             stdlib-only Ollama client (embed + generate/stream)
    rag.py                RAG pipeline: retrieve -> prompt -> generate
    ingest/
      loaders.py          per-format loaders -> text + section metadata
      chunker.py          fixed / sentence / paragraph / recursive splitting
      tags.py             hierarchical tag normalization
      pipeline.py         load->summarize->chunk->embed->store; SUPPORTED_EXTENSIONS
      backfill.py         maintenance CLI: summary backfill, bulk (re-)ingest
    retrieval/
      store.py            the one Qdrant wrapper (vector collection)
      retriever.py        query embedding, top-k + margin + aggregate grouping
      doc_store.py        Qdrant as plain storage: tag taxonomy + conversations
    api/
      app.py              FastAPI wiring only; startup builds the singletons
      services.py         ALL business logic (upload/query/library/delete)
      export.py           answer Markdown -> PDF (mistune + PyMuPDF Story)
      audit.py            append-only JSONL audit log
    ui/                   vanilla JS web app (no build step)
      index.html
      js/                 ES modules: api, chat, library, upload, tags, scope, ...
      styles/             plain CSS: theme, layout, chat, library, tags, ...

  installer/              === THE MANAGER (control panel + CLI) ===
    __main__.py           python -m installer: no args -> panel, args -> CLI
    cli.py                argparse CLI: install/start/stop/status/logs/models/...
    web/server.py         stdlib HTTP server: 127.0.0.1, token-gated, SSE jobs
    backends/             Backend interface + DockerCompose / Portable backends
    installers/           Installer interface + Docker / Portable installers
    model_store.py        filesystem ops on the Ollama store (list/export/import)
    platform_detect.py    OS/arch/GPU detection -> picks the pathway
    process.py            subprocess plumbing; BackgroundJob (thread + queue)
    config.py             app-root/config discovery, config.toml rewriting
    logging_setup.py      logs/installer.log

  deploy/                 === PACKAGING (build machine only) ===
    Dockerfile            builds the app image with uv
    docker-compose.yml    qdrant + ollama + api stack; PINNED versions
    package.sh            builds dist/: bundle tarball + frozen manager + README
    package-macos-portable.sh   vendored CPython + portable runtime for macOS
    data/                 sample documents (also handy in dev)
    ollama_models/        containerized Ollama model store (bind-mounted)
    qdrant_storage/       dev/prod Qdrant data

  tests/
    unit/                 fast, offline: chunker, loaders, retriever, CLI, web, ...
    eval/                 quality suite (pytest -m eval): needs live stack
      corpus_gen.py       generates the synthetic ground-truth corpus
      cases.py            the graded question/answer cases
      test_retrieval.py   precision / recall on retrieval
      test_answers.py     answer correctness, citations, refusals
      test_chunking.py    offline chunk-health checks (run in CI too)
    test_system.py        integration smoke tests (pytest -m integration)

  .vscode/                launch profiles + tasks (Section 4)
\end{lstlisting}

Two files deserve special attention before you change anything:

\begin{itemize}
  \item \textbf{\texttt{schema.py}} is the single shared contract: every
        chunk stored in Qdrant uses exactly those payload fields, and every
        consumer reads them through its helpers. It is what keeps ingestion
        and retrieval from drifting apart. Extend it deliberately; never
        write ad-hoc payload keys.
  \item \textbf{\texttt{config.py}} defines every tunable and its default.
        Anything configurable goes here (and in
        \texttt{config.toml.example}), never as a scattered constant.
        Precedence: environment variable $>$ \texttt{config.toml} $>$
        dataclass default.
\end{itemize}

Layering rules worth preserving: \texttt{app.py} contains HTTP wiring only
--- business logic goes in \texttt{services.py}. The manager's UI layers
(\texttt{web/server.py}, \texttt{cli.py}) talk only to the
\texttt{Backend}/\texttt{Installer} interfaces, never to \texttt{docker
compose} directly --- that is what lets the same code manage both
deployment pathways with no OS conditionals in the UI.

\section{Getting set up (Git + VS Code)}

\subsection{Prerequisites}

\begin{itemize}
  \item \textbf{Git}, \textbf{VS Code} (with the \emph{Python} +
        \emph{Python Debugger} extensions)
  \item \textbf{uv} --- the only Python tooling you need (it
        installs/manages Python 3.14 itself)
  \item \textbf{Docker Engine + Compose v2} --- for Qdrant/Ollama in dev
        and everything in prod. Your user must be able to run
        \texttt{docker} without sudo.
  \item Optional but recommended: an NVIDIA GPU (set \texttt{gpu = true});
        CPU works with the small model.
\end{itemize}

\subsection{Clone and bootstrap}

\begin{lstlisting}
git clone https://github.com/ashwathram/dffrnt-airgapped-ai-assistant.git
cd dffrnt-airgapped-ai-assistant
uv sync                        # creates .venv from uv.lock (incl. dev group)
cp config.toml.example config.toml   # then review: gpu flag, llm_model size
\end{lstlisting}

Pick \texttt{llm\_model} to match your hardware before first start --- the
comments in \texttt{config.toml.example} give VRAM figures per size
(\texttt{qwen3:4b} is CPU-viable; \texttt{14b} wants a 16\,GB GPU;
\texttt{30b-a3b} a 24\,GB GPU).

Branch workflow: \texttt{main} is the default branch and PR target; feature
work happens on branches (e.g.\ \texttt{dev-installer-integration}). Open
VS Code with \texttt{code .} and select the \texttt{.venv} interpreter if
it isn't picked up automatically.

\subsection{Run the dev loop}

\begin{lstlisting}
uv run python -m installer dev   # Qdrant + Ollama containers, health-waited,
                                 # configured models pulled/ensured
uv run dffrnt-api                # API + UI on the host -> http://localhost:8000
\end{lstlisting}

Or just use the VS Code launch profile \textbf{DFFRNT AI Assistant (Dev:
host API)}, which does both (\S\ref{sec:launch}). Code changes to the
Python API require a restart (no auto-reload is configured); UI changes
(\texttt{dffrnt\_assistant/ui/}) are plain static files --- refresh the
browser.

\subsection{Run the tests}

\begin{lstlisting}
uv run pytest                        # unit + offline chunk-health (fast, no stack)
uv run pytest -m eval                # quality suite; needs the dev stack up
uv run pytest -m eval --reset-corpus # ...wiping the KB, re-ingesting tests/eval/corpus
uv run pytest -m integration         # smoke tests against a RUNNING API server
\end{lstlisting}

The default \texttt{addopts} in \texttt{pyproject.toml} excludes
\texttt{integration} and \texttt{eval}, so bare \texttt{pytest} is always
safe offline --- that's also what CI runs. The eval suite grades retrieval
precision/recall, answer correctness, citation coverage, and refusal
behaviour against a synthetic corpus generated by
\texttt{tests/eval/corpus\_gen.py}. Answer tests generate with the local
LLM, so a full pass takes a while.

\begin{gnote}
Eval scores are only comparable against the same model and settings. If you
change \texttt{llm\_model}, \texttt{llm\_num\_ctx}, chunking, or
\texttt{score\_margin}, re-run the suite to establish a fresh baseline
before judging your change.
\end{gnote}

\section{Launch profiles (.vscode/launch.json)}
\label{sec:launch}

\begin{pltable}{p{5.2cm} L}
\rowcolor{tablehead}\thd{Profile} & \thd{What it does}\\
\endfirsthead
\rowcolor{tablehead}\thd{Profile} & \thd{What it does}\\
\endhead
DFFRNT AI Assistant (Dev: host API) & The everyday profile. Pre-launch task
brings up Qdrant + Ollama (via \texttt{python -m installer dev}), then runs
the API on the host under \texttt{debugpy} --- breakpoints anywhere in
\texttt{dffrnt\_assistant/}.\\
DFFRNT AI Assistant (Prod: Docker container) & Rebuilds the app image from
source, then starts the full containerized prod stack via the manager
(\texttt{installer start}). Detached --- stop with \texttt{uv run python -m
installer stop}. Use to verify container behaviour.\\
DFFRNT Control Panel (browser) & Runs the manager's browser panel under the
debugger --- for developing the installer/manager itself. Deliberately has
\textbf{no} pre-launch task: the panel manages the stack, so nothing may
pre-start it.\\
Tests: unit & \texttt{pytest -v} (the offline default suite) under the
debugger.\\
Tests: evaluation (live stack) & The eval suite with the stack auto-started
by the pre-launch task. Skips with instructions if stack/corpus is
missing.\\
Tests: evaluation (reset corpus first) & Same, but
\texttt{--reset-corpus}: wipes the document/summary collections and
re-ingests \texttt{tests/eval/corpus} first. \textbf{Destructive to your
dev KB.}\\
Package: AWS deployment & \texttt{deploy/package.sh TARGET\_SYSTEM=AWS} ---
small online bundle into \texttt{dist/}.\\
Package: Offline deployment & \texttt{deploy/package.sh
TARGET\_SYSTEM=OFFLINE} --- air-gapped bundle (images + model store
included).\\
Package: macOS deployment & OFFLINE bundle +
\texttt{package-macos-portable.sh arm64}: vendored CPython, manager as
source, portable container runtime. The app image must be linux/arm64 for
Apple Silicon.\\
\end{pltable}

The pre-launch task \textbf{Setup: All prerequisites}
(\texttt{.vscode/tasks.json}) runs three tasks in parallel: \texttt{uv
sync}, \texttt{mkdir -p data logs}, and \texttt{uv run python -m installer
dev} (GPU-aware, health-waited, models ensured).

\section{Things future developers should know}

\begin{itemize}
  \item \textbf{The payload schema is load-bearing.} All retrieval
        filtering (tags, filenames, summaries) works off payload fields
        defined in \texttt{schema.py}. Adding a feature that needs new
        per-chunk data means: extend \texttt{PAYLOAD\_FIELDS}, write it in
        \texttt{pipeline.build\_payload}, read it through a
        \texttt{schema.py} helper, and re-ingest.
  \item \textbf{\texttt{llm\_num\_ctx} is not a detail.} Ollama silently
        truncates prompts to its context window; the 4096 default eats RAG
        prompts whole. The shipped 16384 was sized to worst-case retrieval
        + \texttt{llm\_num\_predict}. If you widen retrieval
        (\texttt{top\_k}, aggregates) or lengthen history, re-check the
        budget.
  \item \textbf{The embed model is a one-way door.} Changing
        \texttt{embed\_model} / \texttt{vector\_size} / prefixes
        invalidates every stored vector --- that's why the manager exposes
        LLM swapping but deliberately not embedder swapping. After any such
        change: \texttt{dffrnt-manager reingest} (or \texttt{python -m
        dffrnt\_assistant.ingest.backfill --ingest-missing --force}).
  \item \textbf{Generation settings encode hard-won lessons.}
        \texttt{repeat\_penalty = 1.1} stops the reasoning model looping
        the same thought until timeout; \texttt{num\_predict = 5000} is the
        runaway cap that still leaves room for thinking \emph{and} answer
        (3000 guillotined answers). Change them only with the eval suite in
        hand.
  \item \textbf{\texttt{config.toml} is frozen into bundles.}
        \texttt{deploy/package.sh} ships the repo's \emph{live}
        \texttt{config.toml} --- its state at package time (model choice,
        \texttt{gpu} flag, prompt) is exactly what a deployment runs. Check
        it before packaging.
  \item \textbf{Compose pins versions on purpose.} Qdrant/Ollama image
        tags in \texttt{docker-compose.yml} are pinned; bump them there and
        re-run the eval suite before packaging.
  \item \textbf{The system prompt is a security boundary.} Retrieved
        documents are wrapped in \texttt{<documents>} and the prompt
        instructs the model to treat their contents as data --- the
        prompt-injection defence. Keep that instruction (and the exact
        refusal sentence, which the UI/tests rely on) when editing prompts,
        and keep \texttt{config.toml}'s prompt in sync with
        \texttt{DEFAULT\_SYSTEM\_PROMPT} in \texttt{rag.py}.
  \item \textbf{The API has no authentication.} Anyone who can reach port
        8000 can read, upload, and delete. Fine on a trusted LAN behind the
        air gap; adding auth is an obvious extension point (do it in
        FastAPI middleware, and gate the destructive endpoints first).
  \item \textbf{Long-running work must not block.} In the manager,
        anything that runs a subprocess to completion goes through
        \texttt{process.BackgroundJob} (worker thread + output queue; the
        web panel drains it over SSE, the CLI passes \texttt{print}).
        Follow that pattern for new verbs.
  \item \textbf{The Ollama model store is just files.}
        \texttt{manifests/\ldots/<name>/<tag>} JSON plus content-addressed
        \texttt{blobs/sha256-*}. \texttt{model\_store.py} manipulates it
        directly (list/export/import work with the stack down); only
        pull/delete talk to the running Ollama.
  \item \textbf{Logs to know:} \texttt{logs/installer.log} (manager
        operations), \texttt{logs/audit.jsonl} (app audit trail),
        \texttt{dffrnt-manager logs <svc>} (containers).
\end{itemize}

\section{Extension recipes}

\paragraph{Add a document format.}
Write a loader in \texttt{ingest/loaders.py} returning the standard
Document dict (text + sections with page/heading metadata), register the
extension in \texttt{load\_file}'s dispatch and in
\texttt{SUPPORTED\_EXTENSIONS} (\texttt{ingest/pipeline.py}), add a unit
test with a fixture file, and (optionally) surface the extension in the UI
upload validator (\texttt{ui/js/upload.js}).

\paragraph{Add an API endpoint.}
Put the logic in \texttt{api/services.py} (raise \texttt{UserError(status,
detail)} for client errors), add the route in \texttt{api/app.py} as a thin
adapter, wire the UI call in \texttt{ui/js/api.js}, and consider whether
the action belongs in the audit log.

\paragraph{Add a manager verb.}
Implement it on the \texttt{Backend} interface (both pathways if
applicable), expose it in \texttt{cli.py} (argparse) and
\texttt{web/server.py} (plus the panel UI in
\texttt{installer/web/static/}), and run it through
\texttt{BackgroundJob} if it can take more than a moment.

\paragraph{Tune retrieval quality.}
The knobs are \texttt{top\_k}, \texttt{score\_margin},
\texttt{aggregate\_group\_limit/size}, and the chunking block. The workflow
is: change one thing $\rightarrow$ \texttt{pytest -m eval --reset-corpus}
$\rightarrow$ compare. Never tune by vibes; the margin cut especially is
corpus-dependent.

\section{FAQ --- design choices and trade-offs}

\paragraph{Why no LangChain / LlamaIndex / RAG framework?}
The pipeline is $\sim$a dozen small modules; a framework would add hundreds
of transitive dependencies to an air-gapped bundle, obscure the retrieval
logic this project exists to tune, and churn APIs faster than a
student-team handover cycle. Trade-off: we reimplement loaders and prompt
assembly ourselves --- accepted, because those are exactly the parts we
want control over (and they're covered by unit + eval tests).

\paragraph{Why a stdlib-only Ollama client (\texttt{urllib})?}
One fewer dependency to vendor, and the Ollama HTTP API is small (two
endpoints used). Trade-off: no connection pooling or typed responses;
mitigated by the client being $\sim$200 lines with tests.

\paragraph{Why Ollama rather than vLLM / llama.cpp directly?}
Ollama gives model lifecycle management (pull, quantized model store,
keep-alive), a stable HTTP API, and one identical serving story on Linux
containers \emph{and} native macOS with Metal. Trade-off: less throughput
tuning than vLLM --- irrelevant at single-office concurrency.

\paragraph{Why Qdrant?}
Runs as a single small container fully offline, has payload filtering (tag
scoping needs it), grouped search (the aggregate-query feature), and
persistent named collections. Trade-off: heavier than an embedded library
(e.g.\ sqlite-vec), but it's also reused as the app's only datastore ---
see next question.

\paragraph{Why are tags and conversations stored in Qdrant with dummy vectors?}
\texttt{doc\_store.py} attaches a throwaway 1-D vector because Qdrant
requires one, then only ever upserts/fetches by id. This keeps the
deployment at \textbf{zero additional databases} --- no SQLite file to
migrate, no second container. Trade-off: mildly unidiomatic; documented in
the module docstring.

\paragraph{Why a vanilla-JS UI with no build step?}
An air-gapped target can't \texttt{npm install}, and a build pipeline is
one more thing a future team must resurrect before touching the UI. Native
ES modules + plain CSS mean ``edit file, refresh browser'' forever.
Trade-off: no reactivity framework --- state is managed manually in
\texttt{js/state.js}, which is fine at this UI's size but would strain if
the UI grew much larger.

\paragraph{Why is the model qwen3 and the embedder bge-m3?}
\texttt{qwen3} ships in sizes spanning CPU-only to 24\,GB GPU with strong
multilingual + reasoning behaviour under Ollama; \texttt{bge-m3} is a
robust multilingual embedder with 1024-dim vectors and no task prefixes.
Both are swappable in config --- but swapping the embedder invalidates the
index (Section~5), and eval baselines are per-model.

\paragraph{Why does the manager render in a browser instead of Tk/Qt?}
The browser is the one GUI runtime every target already has. A Tk build
needs per-platform freezing and made macOS support painful; a localhost web
panel is identical everywhere, works over SSH tunnels, and shares its
plumbing with the headless CLI. Trade-off: a local HTTP server needs a
security model --- hence loopback-only binding + per-session token.

\paragraph{Why is dffrnt-manager a PyInstaller binary on Linux/Windows but source + vendored CPython on macOS?}
Targets must not need Python. PyInstaller freezes per-platform, and the
team had no Mac in the build loop --- so the macOS bundle carries the
manager as source plus a relocatable CPython (python-build-standalone) and
a shell launcher. Same UX (\texttt{./dffrnt-manager}), no Mac needed to
build.

\paragraph{Why Python $\geq$ 3.14?}
Recent stdlib carries weight here: \texttt{tomllib} (config), modern
typing, and current security fixes for a long-lived offline install. uv
makes the interpreter version a non-issue for developers.

\paragraph{Why FastAPI?}
Streaming responses (ndjson token streams), static file mounts for the UI,
pydantic validation at the edge, and wide familiarity for onboarding teams.
The app deliberately keeps it to a thin adapter so a framework swap would
be contained to \texttt{app.py}.

\paragraph{Why is there a separate summaries collection?}
Chunk-level search alone misses ``about-ness'': aggregate questions
(``compare all proposals'') need to know \emph{which documents} are
relevant, not just which 512-char chunks. Per-document summaries give a
routing signal and prompt briefing. Trade-off: summaries are LLM-generated
at ingest time (slower ingestion) and must be backfilled for pre-existing
docs (\texttt{ingest/backfill.py}).

\paragraph{Why does bare pytest skip the interesting tests?}
So CI and any offline checkout stay green with no stack. The quality suite
is opt-in (\texttt{-m eval}) because it needs Qdrant + Ollama and real
generation time. The skip messages tell you exactly how to bring the stack
up.

\paragraph{Where would authentication go if we added it?}
FastAPI dependency/middleware in \texttt{app.py} gating the mutating routes
first (upload/delete/tags), then the query routes. The manager panel
already has its token model and needs nothing. Remember the UI's
\texttt{js/api.js} is the single place all fetches go through --- one spot
to attach credentials.

\end{document}
